% !TeX root = ../main.tex
\chapter{Compact Students and Embedded Inference Design}
\label{ch:student}
\section{Why a Random-Forest Student Is Selected}
The student converts the teacher-generated examples into a policy whose inputs
and implementation suit the intended onboard task. Its input is one numeric
LiDAR vector and its output is one of the observed steering labels. A random
forest can express nonlinear combinations of range thresholds without using
the teacher's map geometry or vehicle-state features~\cite{breiman2001random}.
It is the principal student family selected for this project.

There are three practical reasons for this choice. First, threshold-based
partitions directly operate on the binned distance representation. Second,
maximum depth and tree count bound the number of decisions traversed during
prediction, exposing a capacity--computation trade-off. Third, fitted tests and
leaf outputs can be exported into ordinary C++ inference. Only prediction is
intended for the board; tree construction and RL optimisation stay on the host.
The deployment path therefore avoids carrying the Python training environment
onto the vehicle.

The ensemble combines multiple tree predictions, providing a motivation for
using a forest instead of trusting one partition of the demonstrations~\cite{breiman2001random}.
However, more trees also require additional stored structure and aggregation.
Neither ensembling nor dropout augmentation guarantees robustness to missing
angular information. The closed-loop experiments test those assumptions rather
than treating them as consequences of the model name.

A compact or quantised neural student is also a possible design, and a deeper
single tree may be cheaper than a forest. The present study does not contain
a board-level comparison of all these alternatives. It includes a single-tree
baseline to test whether a simpler decision structure retains useful behaviour.
Selecting a forest is thus a motivated starting point for the resource-limited
objective, not a requirement that the final evidence must favour forests.

\section{Teensy 4.0 Hardware Budget}
The intended target is the Teensy 4.0. Table~\ref{tab:teensy_resources} gives
the manufacturer's nominal resources~\cite{pjrcTeensy40}. Flash is the
non-volatile program store relevant to the model; calling it simply ``ROM''
would obscure the distinction between rewritable firmware storage and boot ROM.

\begin{table}[htbp]
\centering\small
\caption{Nominal Teensy 4.0 resources from PJRC documentation.}
\label{tab:teensy_resources}
\begin{tabularx}{\textwidth}{p{2.5cm}>{\raggedright\arraybackslash}X}
\toprule Resource & Specification \\
\midrule
CPU & Arm Cortex-M7, 600~MHz; hardware single- and double-precision floating point. \\
RAM & 1,024~KiB total: 512~KiB RAM1 (tightly coupled) and 512~KiB RAM2 (DMA-oriented). \\
Flash & 2~MiB physical; upper 64~KiB reserved, leaving 1,984~KiB for the program region. \\
\bottomrule
\end{tabularx}
\end{table}

The RAM total is not a single unrestricted model buffer. Under the documented
default placement, functions are copied from Flash to ITCM in RAM1; stack and
ordinary variables also consume RAM1. RAM2 serves larger buffers and the heap.
Consequently, generated inference code can consume both Flash and RAM1, depending
on compiler and linker placement~\cite{pjrcTeensy40}. A linker map is needed to
determine the actual allocation.

This budget explains the training/deployment split. The board need only retain
the selected inference implementation and its working data, rather than the
Gazebo environment, PPO training state or demonstration dataset. CPU frequency
alone does not establish a decision latency, and fitting within nominal memory
does not establish low power. Those questions require the integrated firmware
and measurements under its sensor and motor-control workload.

\section{Resource Reasoning and Deployment Criteria}
For $T$ binary trees of maximum depth $D$, an inference path evaluates at most
$TD$ internal-node thresholds, excluding input processing and leaf aggregation.
A full binary tree at that depth has at most $2^{D+1}-1$ nodes. These structural
bounds make the configured trade-off explicit: an 88-tree depth-three forest
uses at most 264 threshold comparisons and 1,320 total nodes, whereas a single
depth-six tree uses at most six comparisons and 127 nodes. Actual trees can
contain fewer nodes. These are derived bounds, not measurements of execution
time or memory on the board.

The exported forest also accumulates leaf probabilities for the six fitted
classes before selecting an action. That work, stored thresholds, leaf values,
compiled instructions and sensor buffers all contribute to deployment cost.
Consequently, a shallow forest is not automatically smaller or lower-power than
every neural or single-tree alternative. Header size and host timings are
useful development indicators, but cannot establish the final resource budget.

Inspection of the frozen fitted artifacts gives the exact structures in
Table~\ref{tab:artifact_resources}. All three forests contain 1,320 nodes and
704 leaves; T6 contains 127 nodes and 64 leaves. Thus the full-tree bounds are
attained by these particular artifacts. The headers express these structures
as branches and constants rather than serialising the Python estimator onto
the board. Their textual size is affected by numeric formatting and generated
syntax, whereas compiled storage depends on instruction generation, constant
placement and optimisation.

\begin{table}[htbp]
\centering\small
\caption[Frozen model structures and file sizes]{Frozen model structures and file sizes, checked against
the model hashes used in the closed-loop reports. File sizes are bytes, not
compiled Flash or runtime RAM.}
\label{tab:artifact_resources}
\begin{tabular}{lrrrrr}
\toprule Model & Trees & Nodes & Leaves & Joblib bytes & C++ source bytes \\
\midrule
RF-D & 88 & 1,320 & 704 & 187,946 & 214,856 \\
RF-B & 88 & 1,320 & 704 & 187,943 & 218,429 \\
RF-R & 88 & 1,320 & 704 & 187,942 & 196,147 \\
T6 & 1 & 127 & 64 & 16,727 & 21,178 \\
\bottomrule
\end{tabular}
\end{table}

The exported interface takes a pointer to 200 single-precision inputs and
uses six double-precision probability accumulators. With four-byte
\texttt{float} and eight-byte \texttt{double}, these arrays correspond to
800 and 48 bytes respectively. Their combined 848 bytes is a source-level
accounting of two arrays, not a measured peak RAM requirement. Raw LiDAR
buffers, feature conversion, stack frames, interrupts, runtime libraries and
code copied into RAM remain outside that count. The existing export preserves
double-precision accumulation; reducing that precision would be a model-export
change requiring a new agreement check.

The structural cost comparison is also narrower than a timing comparison.
The actor's dense layers contain 55,680 weights, each participating in a
multiply--accumulate in a direct dense forward pass, plus biases and nonlinear
activations. The forests traverse at most 264 thresholds and aggregate leaf
probabilities. These operation types have different costs and memory-access
patterns, so their counts cannot be divided to claim an execution-speed ratio.
They explain why a compact student is worth testing, while leaving the board
benchmark as the necessary final check.

The intended acceptance criteria are that compiled code and constants fit in
available Flash, peak runtime storage fits in available RAM, and sensing,
preprocessing, inference and command publication meet the selected control
period. Energy per decision and average board power must be measured under the
actual workload to substantiate the low-power objective. The simulation decision
interval is 0.2~s; a physical implementation preserving that interface would
need to meet the whole interval, not merely an isolated classifier benchmark.
These criteria guide later integration and are not presented as completed tests.

\section{Data Validation and Map Separation}
The student workflow validates the distance and label fields before fitting.
Non-finite values, out-of-range distances, unsupported labels and all-zero
scans are rejected according to the configured checks. Zero values in active
bins are otherwise permitted, since they represent a condition explicitly
included in the training data. When original and augmented files are used
together, row counts and labels are checked for alignment.

Keeping related observations together follows the rationale of grouped
evaluation for dependent samples~\cite{sklearnCV}. Here a fixed map-based
holdout is used, rather than a k-fold cross-validation procedure.
Map grouping determines the fitting/validation split. T3, R1 and R3 provide
1,079,984 original rows, with the same number of dropout copies. E2 provides
359,999 original validation rows and equally many copies. Thus the combined
variant uses 2,159,968 fitting rows and 719,998 validation rows. The saved
models are not subsequently refitted using E2. This preserves the reported
fitting exclusion, although repeated use of E2 for model choices makes it a
validation set rather than a final test set.

All main models retain the 200-column representation. Feature selection is
available in the pipeline but is not used in this comparison. Keeping fixed
zero columns maintains the input contract; it does not add useful angular
measurements. Feature selection results should therefore not be claimed as
a contribution of the reported experiments.

\section{Model Configurations}
The students use scikit-learn's classifier implementation~\cite{pedregosa2011sklearn,sklearnForest}.
The main forests use 88 trees, maximum depth three, minimum leaf size one,
no class weighting and seed 42. Three data variants are compared: original
only (RF-R), dropout only (RF-D), and original plus dropout (RF-B). The single
tree (T6) uses the combined data and a maximum depth of six. It is implemented
through the same classifier/export interface with one estimator, bootstrap
disabled and all features eligible at a split.

\begin{table}[htbp]
\centering\small
\caption[Student model configurations]{Main student configurations. RF settings use the classifier defaults
for bootstrap and feature sampling.}
\label{tab:models}
\begin{tabular}{lrrrr}
\toprule & RF-D & RF-B & RF-R & T6 \\\midrule
Trees & 88 & 88 & 88 & 1 \\
Maximum depth & 3 & 3 & 3 & 6 \\
Training variant & Dropout & Both & Raw & Both \\
Bootstrap & Yes & Yes & Yes & No \\
Split features & $\sqrt{200}$ & $\sqrt{200}$ & $\sqrt{200}$ & All \\
Input columns & 200 & 200 & 200 & 200 \\
\bottomrule
\end{tabular}
\end{table}

The single-tree depth was compared at 3, 6, 8 and 10 using the validation
results. Depth six was retained; the corresponding combined-input accuracies
were approximately 0.6043, 0.6653, 0.6619 and 0.6044. This selection is part of
the development process. It does not establish that depth six would be optimal
on a new map or under a new sensor distribution.

Table~\ref{tab:models} exposes an important interpretation limit. The single
tree and forests differ in depth, number of estimators, bootstrap and split
feature sampling. Their results compare these complete configurations, not an
isolated causal effect of ensembling. Similarly, the combined data variant
contains twice as many rows as either individual variant, so it does not
separate augmentation content from all sample-count effects.

\section{Prediction and Offline Metrics}
For tree $j$, let $p_j(k\mid x)$ be the class distribution at the leaf reached
by scan $x$. The deployed prediction is
\begin{equation}
 \bar p(k\mid x)=\frac{1}{T}\sum_{j=1}^{T}p_j(k\mid x),\qquad
 k^*=\arg\max_k \bar p(k\mid x),\qquad \hat a(x)=c_{k^*}.
\end{equation}
Probability averaging follows the classifier's documented prediction
rule~\cite{sklearnForest}; $c_k$ maps an internal class index to the original steering command.
This mapping is necessary because the observed labels are not a contiguous
sequence beginning at zero. The C++ exporter retains leaf probabilities and
class ordering from the fitted estimator.

Exact action accuracy measures agreement with the teacher. Steering mean
absolute error (MAE) and one-step tolerance accuracy are
\begin{equation}
 \mathrm{MAE}=\frac{1}{N}\sum_i|\hat a_i-a_i|,\qquad
 A_{\pm1}=\frac{1}{N}\sum_i\mathbf{1}\{|\hat a_i-a_i|\leq1\}.
\end{equation}
The label distance is meaningful here because the wheel-speed mapping varies
linearly with steering label. Nevertheless, these are imitation metrics,
not direct measures of collision risk.

Class imbalance motivates macro F1 and per-class support. Macro F1 is the
unweighted mean of per-class F1 scores; the selected labels and handling of
undefined scores affect the result~\cite{sklearnF1}. Two explicit
conventions are retained: F1 averaged over labels occurring in the evaluated
targets, and F1 averaged over the fixed set $\{5,6,7,8,9,12\}$ with undefined
class scores set to zero. These values can differ even for the same prediction
vector. In particular, the later E2 batch contains no label 12. Tables identify
the convention rather than treating all reported macro F1 values as comparable.

A majority-class accuracy of approximately 0.6241 on the original E2
validation distribution provides a simple offline reference. It is not used
as an autonomous driving baseline: continuously selecting the majority
straight label does not evaluate the same navigation capability as a trained
feedback policy.

\section{Host-Side C++ Export Verification}
The exporter writes nested threshold tests and accumulates leaf class
probabilities before selecting a command. Checks include sampled observations,
random probes and cases around split thresholds. The three primary forests
each pass 3,149 verification inputs, while T6 passes 1,490. These counts refer
to tested inputs and should not be interpreted as a formal proof over every
possible input vector. Model artifacts and reports are identified in
Appendix~\ref{sec:evidence_models}.

Publication of a model follows successful validation so that an unsuccessful
export does not silently replace an earlier usable artifact. Reports retain
model configuration, software versions and header fingerprints. This provides
a link between fitted predictions and the model executed in subsequent
experiments.

The depth-six tree produces a substantially smaller source header than the
88-tree forests in the archived artifacts, and its recorded host Python
inference is faster. These are useful development observations, but neither
source-file bytes nor Python timings are target-board resource measurements.
No claim about compiled Flash, working RAM or complete embedded control-period
compliance is made from them.

\section{Frozen Candidates and Cross-Batch Use}
The cross-batch study evaluates the four existing models on both collections
without retraining. It retains map-level, episode-level and per-class results.
The 64 map-level cells arise from two batches, four models, two input variants
and four maps. The 1,120 episode-level metric rows arise from re-evaluating the
80 and 60 teacher episodes under each model/input combination; they are not
1,120 newly driven student episodes.

T6 and RF-B are retained as the two frozen candidates for the proposed
independent-map experiment. RF-R and RF-D remain important ablation models.
Freezing this candidate set does not declare a universal winner, and the
proposed experiment has not yet supplied evidence. Keeping the later batch
out of fitting allows its results to describe the original models rather
than a moving sequence of retrained controllers.
